\section{The Last Cracker}
\label{sec:n-last-cracker}
\source{ICPC 2014--15 Kharagpur Online Re-contest (PYQ, ACM14KG2)}{Easy}

\problemstatement{$N$ crackers stand in a circle. The spark starts at
cracker $a$; each cracker fires and passes the spark to the next ($N\to1$).
The spark dies during pass $M+1$. Which cracker fires last?
($N,M\le1000$.)}

\begin{patternbox}
Circular indexing $\Rightarrow$ convert to $0$-based, add, take $\bmod N$,
convert back.
\end{patternbox}

\subsection*{Approach and intuition}

After $M$ successful passes the spark is at the cracker $M$ steps after $a$;
that cracker fires and the next pass fails. In $0$-based terms the position is
$(a-1+M)\bmod N$, so the answer is $(a-1+M)\bmod N+1$.

\subsection*{Algorithm}
\begin{algorithmic}[1]
\State \Return $(a-1+M)\bmod N+1$
\end{algorithmic}

\dryrun{$N=5$, $a=1$: $M=3\to4$, $M=5\to1$, $M=19\to(19\bmod5)+1=5$. \checkmark}

\subsection*{C++17 implementation}
\cppfile{number-theory/02_the_last_cracker.cpp}

\begin{complexitybox}
\textbf{Time:} $O(1)$. \textbf{Space:} $O(1)$.
\end{complexitybox}

\begin{pitfallbox}
\begin{itemize}
  \item $(a+M)\bmod N$ gives $0$ instead of $N$ --- always shift to $0$-based
        first.
  \item $M=0$: the starting cracker fires and the spark dies; answer $a$.
\end{itemize}
\end{pitfallbox}
